\subsection{What the guarantee covers}
\label{sec:guarantee}\label{sec:assumptions}

\Cref{prop:simultaneous} is a conditional statement: if the coefficient region
covers the true coefficients with probability $1-\alpha$, the certificate
inherits that probability for every generator and every tolerance. Everything
else is in the condition. \Cref{tab:guarantees} lists each guarantee in the
paper with what it requires, and the rest of this subsection takes the
requirements in turn.

\begin{table}[!htbp]
\centering\small
\begin{tabular}{>{\raggedright\arraybackslash}p{0.26\textwidth}>{\raggedright\arraybackslash}p{0.39\textwidth}>{\raggedright\arraybackslash}p{0.24\textwidth}}
\toprule
guarantee & requires & evidence \\
\midrule
False-certification rate $\le \alpha$, simultaneously over generators and
tolerances (\Cref{prop:simultaneous}) &
a region with coverage $\ge 1-\alpha$ for coefficients of a class that
contains the truth; dictionary, candidate class, target measure $\nu$ and
metric $W$ fixed without looking at the data &
\cref{tab:calibration}, \cref{fig:headline}a \\
\addlinespace
Exact finite-sample coverage of the region &
Gaussian homoskedastic errors on the right-hand side, independent of states
that are observed exactly &
\Cref{tab:robustness} \\
\addlinespace
Asymptotic coverage of the region &
independent rows (White sandwich), independent trajectories with their number
growing (cluster-robust), or long trajectories (Newey--West) &
\Cref{tab:robustness}, \SFdependence \\
\addlinespace
Validity when the class is chosen from the data
(\Cref{prop:classes}) &
the largest class of a nested family contains the truth; regions of
\Cref{prop:closedform}(iii) (Gaussian homoskedastic errors, fixed design,
exactly observed states) &
\Cref{sec:selection} \\
\addlinespace
Validity under a misspecified class &
an externally justified bound on the omitted part (sensitivity parameter) &
\SFmodelerror \\
\addlinespace
Coverage diagnostic $\kappa$ (\Cref{prop:kappa}) & none to compute it from
the observed states; as a bound on the true defect, the fitted class contains
the truth (a population-level statement; estimation error not included) &
\cref{fig:headline}b, \STtrajectory \\
\addlinespace
Resolution limit (\Cref{prop:lecam}) & Gaussian errors, fixed design, one
generator, fixed model class & \cref{fig:headline}c, \STqsweep \\
\midrule
\emph{Not covered} & noisy states (errors in variables); a class chosen by a
lack-of-fit test on the same data; translation generators; discrete
symmetries & \SFeiv, \Cref{sec:selection} \\
\bottomrule
\end{tabular}
\caption{\textbf{Guarantees and the assumptions they rest on.} Each row states
what has to hold for the guarantee in the first column; the last column points
to the experiments that probe it. The first row is the main result; all others
qualify it.}
\label{tab:guarantees}
\end{table}

\paragraph{The model class.} Validity needs a coefficient vector $\beta^\star$
with $F = F_{\beta^\star}$, so the fitted class must contain the truth. Three
situations have to be kept apart.

\emph{(i) A class fixed in advance and correctly specified.} This is the
setting of \Cref{prop:simultaneous}, and the false-certification rate is at
most $\alpha$.

\emph{(ii) A class selected from the data.} A natural workflow enlarges the
class until a lack-of-fit $F$-test against a larger class stops rejecting, and
then certifies in the chosen class on the same data. This has no guarantee, for
two reasons. A non-rejection is not evidence that the class contains the truth
(the test may simply lack power; this is the same inversion of the burden of
proof that motivates the equivalence formulation), and the region in a
data-selected class is not a post-selection confidence region. The first
reason can be quantified.

\begin{proposition}[What a lack-of-fit test cannot see]
\label{prop:lof}
Let the truth be $F = F_{\beta^\circ} + h$, where $F_{\beta^\circ}$ is the
least-squares projection of $F$ on the fitted class over the design and $h$
lies in a larger class $\mathcal H$ of $m_{\mathcal H}$ functions per
component. The $F$-test of the fitted class ($m$ functions per component)
against $\mathcal H$ has noncentrality $\lambda = N\|h\|^2_{\mu_N}/\sigma^2$,
where $\|\cdot\|_{\mu_N}$ is the root-mean-square over the design, and power
$\gamma(\lambda)$ given by the noncentral $F$ distribution with $n(m_{\mathcal
H}-m)$ and $n(N-m_{\mathcal H})$ degrees of freedom. Let $\lambda_\gamma$ solve
$\gamma(\lambda_\gamma)=\gamma$ and
\begin{equation}
  G(\xi) \;=\; \max_{h \in \mathcal H_\perp}
  \frac{\|[\xi,h]\|_{\nu,W}}{\|h\|_{\mu_N}},
  \label{eq:invisible-gain}
\end{equation}
with $\mathcal H_\perp$ the fields of $\mathcal H$ orthogonal to the fitted
class over the design. Then every omitted part with
$\|h\|_{\mu_N} \le \sigma(\lambda_\gamma/N)^{1/2}$ is detected with probability
at most $\gamma$. Each such $h$ adds to the bracket $[\xi,F_{\beta^\circ}]$ a
term with $\|[\xi,h]\|_{\nu,W} \le G(\xi)\,\sigma(\lambda_\gamma/N)^{1/2}$, so it
changes the numerator of the defect of $\xi$ by at most this amount. One
such part, $h^\star$, attains the bound on $\|[\xi,h]\|_{\nu,W}$; if $\xi$ is an
exact symmetry of $F_{\beta^\circ}$, the numerator then changes by exactly this
amount, and the truth $F = F_{\beta^\circ} + h^\star$ has
\begin{equation}
  \defect_\nu(\xi;F) \;=\; \frac{G(\xi)\,\sigma\,(\lambda_\gamma/N)^{1/2}}
  {\bigl(\|DF\,\xi\|^2_{\nu,W} + \|D\xi\,F\|^2_{\nu,W}\bigr)^{1/2}}
  \label{eq:invisible}
\end{equation}
while the test of the fitted class rejects with probability at most $\gamma$.
\end{proposition}

To first order in $h^\star$ the denominator of~\eqref{eq:invisible} is that of
$F_{\beta^\circ}$. The invisible defect therefore scales like $\sigma/\sqrt N$, as the
width of the certificate does, so whether a lack-of-fit gate protects the
certificate is decided by constants: by the derivative amplification in
$G(\xi)$, which grows with the degree of the omitted terms, and by the degrees
of freedom of the test. \Cref{sec:selection} shows both regimes: one where the
test is the more sensitive instrument, and one where misspecification that it
misses with probability $0.5$ can produce true defects of about $1.7$ times the
tolerance and the gate certifies falsely in nearly half of the replications. A
valid procedure for case (ii) needs a class that is assumed to contain the
truth. Given that assumption, choosing among its nested subclasses costs only a
split of $\alpha$ (\Cref{prop:classes}); the assumption itself costs power,
because a larger class widens the region (\Cref{sec:selection}).

\begin{proposition}[Certifying simultaneously over a nested family]
\label{prop:classes}
Under model~\eqref{eq:model} with a fixed design, let $\mathcal F_1 \subset
\dots \subset \mathcal F_K$ be nested linear model classes and assume $F \in
\mathcal F_K$. Give each class the region $\mathcal B_k$ of
\Cref{prop:closedform}(iii) at level $\alpha/K$, computed in that class. For $k < K$ write $h_k(\beta)$ for the part
of $F_\beta \in \mathcal F_K$ omitted by the least-squares projection on
$\mathcal F_k$ over the design, let $\bar\eta_k = \sup_{\beta \in \mathcal B_K}
\|h_k(\beta)\|_{\nu,W}$, and let $U_k$ be the certificate computed in
$\mathcal F_k$ with the $L^2$ model-error term $\bar\eta_k$ (and $U_K$ the
certificate in $\mathcal F_K$). Then with probability at least $1-\alpha$,
$\defect_\nu(\xi_\theta;F) \le \min_k U_k(\theta)$ for every $\theta$.
\end{proposition}

The reported bound is the smallest over the family. Each $\bar\eta_k$ is one
trust-region maximum, because $h_k(\beta)$ is linear in $\beta$. Sample
splitting, choosing the class on one half and certifying on the other, removes
the reuse of the data but does not make a chosen class contain the truth.
Certifying in the largest class alone, at level $\alpha$, is valid under the
same assumption as \Cref{prop:classes}. The proposition's bound is never larger
than that class's certificate at level $\alpha/K$; when the smaller classes do
not help, it is slightly wider than certifying in the largest class at level
$\alpha$, because it still pays for the split of $\alpha$.

\emph{(iii) Misspecification bounded from outside.} When no available class
contains the truth, replace $F$ by $F_\beta + h$ and pay for $h$ explicitly. In
\emph{supremum form}, with $W = I$, if $\sup_{\operatorname{supp}\nu} \|h\| \le
\eta_0$ and $\sup_{\operatorname{supp}\nu}\|Dh\|_F \le \eta_1$, then $[\xi,h]$
moves the numerator by at most $\eta_1\|\xi\|_{\nu} + \eta_0\|D\xi\|_{\nu}$,
where $\|\xi\|_\nu = \|\xi\|_{\nu,I}$ and $\|D\xi\|_\nu = (\int
\|D\xi\|_F^2\,d\nu)^{1/2}$, and each denominator block by at most one of those
terms; valid for any smooth $h$ and conservative, because it spends supremum
norms. In \emph{$L^2$ form}, if the truth lies in a larger class and the
omitted part has $\|h\|_{\nu,W} \le \eta_2$, the numerator moves by at most
$\|\Lambda_\theta\|_{\mathrm{op}}\eta_2$, where $\|\Lambda_\theta\|_{\mathrm{op}}
= \sup_h \|[\xi_\theta,h]\|_{\nu,W}/\|h\|_{\nu,W}$ over $h$ in the larger class
is the largest gain of the linear map $h \mapsto [\xi_\theta,h]$, computed
exactly as a matrix norm; the two denominator blocks are treated likewise. Both forms are expensive
because the Lie derivative \emph{differentiates}: a model error that is small in
size need not be small in slope, and the conversion constant grows with the
degree of the class the error is allowed to occupy. None of the bounds
$\eta_0$, $\eta_1$ and $\eta_2$ can be verified from the data alone, so we
treat them as sensitivity parameters and report the certified tolerance as a
function of them, in the manner standard for
unverifiable assumptions \citep{rosenbaum2002observational}.

\paragraph{The errors.} Exact coverage of the least-squares region needs
Gaussian homoskedastic errors on the right-hand side, independent of states
that are themselves observed without error. Any other construction of
$\mathcal B$ can be supplied instead, and the certificate is then exactly as
valid as that construction: White's sandwich \citep{white1980heteroskedasticity}
for independent rows with arbitrary heteroskedasticity, asymptotically in the
number of rows; the cluster-robust sandwich for independent trajectories with
arbitrary dependence within each, asymptotically in the number of trajectories,
which must exceed the number of coefficients for the region to be bounded in
every direction, and in our experiments exceed it several times over before the
coverage approaches its nominal level; the Newey--West sandwich \citep{newey1987simple} with an
automatic bandwidth \citep{andrews1991heteroskedasticity} for serially
correlated errors along long trajectories. Noisy \emph{states} are a different
matter: the regressors then carry error, least squares estimates a
pseudo-true field, and no covariance estimate repairs the bias. That case is
not covered (\Cref{sec:robustness}).

\paragraph{The target domain, the metric and the candidate class.} The claim
is about the user's $\nu$ and $W$, and nothing is asserted outside $\nu$;
coverage of $\nu$ by the data enters through the design and is quantified by
\Cref{prop:kappa}. We find connected subgroups whose generators lie in
$\mathcal V$, excluding translations, for which the relative defect
degenerates (\Cref{sec:setup}). Discrete symmetries are outside this
framework, as they are outside the Lie-algebraic framework generally.

\paragraph{Learned features.} The construction applies verbatim to any model
that is linear in its parameters given its features, including the last layer
of a network. The classical region is valid only if the features are fixed
independently of the data used for the region, either chosen in advance (as
our polynomial and dictionary features are) or learned on separate data, with
sample splitting or cross-fitting \citep{chernozhukov2018double}. Features
learned on the same data need a region that accounts for that learning.
